
\section{1. Introduction}

\subsection{1.1 Motivation}

Transformer attention at inference has two distinct regimes that pull in 
opposite directions:

\begin{itemize}
\item \textbf{Local regime.} Attention concentrates on nearby tokens, especially in
\end{itemize}
  early layers. Sliding-window attention (Beltagy et al., 2020) exploits this 
  with $\mathcal{O}(L w)$ complexity, where $w$ is the window size.

\begin{itemize}
\item \textbf{Global regime.} Some attention heads (e.g., attention sinks, coreference
\end{itemize}
  heads) attend to distant tokens. Linear attention (Katharopoulos et al., 
  2020; Choromanski et al., 2021) approximates global attention in 
  $\mathcal{O}(L)$ via a kernel factorization.

\textbf{The problem with picking one.} Longformer's window truncates global 
context. Performer's random features degrade local precision. Neither 
extremes is universally better.

\textbf{The opportunity.} Rather than commit to one regime, \textbf{hybrid attention} 
runs both and combines them. If the combination weight is **learned per 
token**, the model can allocate local precision where it matters and global 
approximation elsewhere.

\subsection{1.2 Problem Statement}

Given a pretrained transformer with dense attention, we seek a modified 
module $\mathcal{A}'$ such that:

1. \textbf{Parallel streams:} Both local and global attention are computed 
   simultaneously.
2. \textbf{Learned fusion:} A per-token gate mixes the two streams.
3. \textbf{Stability:} The gate is initialized to favor the pretrained behavior.
4. \textbf{Quality:} $\Delta\text{Acc} \leq {{ACC_TOLERANCE}}\%$ on SST-2.
5. \textbf{Pretrained compatibility:} No base retraining; only LoRA + gate 
   parameters are trainable.

\subsection{1.3 Our Approach: NHA}

\textbf{Native Hybrid Attention} (NHA) runs two attention streams in parallel:

1. \textbf{Local stream} $\mathcal{A}_\text{local}$: sliding-window attention with 
   window size $w$ and chunked computation.
2. \textbf{Global stream} $\mathcal{A}_\text{global}$: RFF-based linear attention 
   with orthogonal projection matrix.

The streams are combined via a \textbf{token-level learned gate}:

\[
O_i = g_i \cdot O_i^\text{local} + (1 - g_i) \cdot O_i^\text{global}
\tag{1.1}
\]

where $g_i = \sigma(w^\top x_i + b)$ with $b$ initialized to $3.0$ (biasing 
toward local, $\sigma(3) \approx 0.95$).

\textbf{Key design choices:}

1. \textbf{Orthogonal RFF projection.} We build the RFF projection via block QR, 
   reducing variance compared to i.i.d. Gaussian projections.
2. \textbf{Chunked causal prefix scan.} For decoders, the RFF stream maintains 
   cumulative state across chunks, preserving strict causality.
3. \textbf{Token-level gating.} The gate is per token, not per head, capturing 
   position-specific needs without excess parameters.

\subsection{1.4 Contributions}

\begin{itemize}
\item \textbf{A native hybrid attention module.} NHA runs local and global streams in
\end{itemize}
  parallel within each layer, fused by a learned token-level gate.
\begin{itemize}
\item \textbf{Orthogonal RFF.} We use block-QR orthogonal projections for reduced
\end{itemize}
  variance.
\begin{itemize}
\item \textbf{Chunked causal scan.} A memory-bounded RFF implementation for decoders.
\end{itemize}
\begin{itemize}
\item \textbf{A production-hardened implementation.} Supports four architectures
\end{itemize}
  with RoPE, GQA, LoRA recovery, and safe injection.
\begin{itemize}
\item \textbf{A systematic evaluation.} Benchmarked on DistilBERT, RoBERTa, and GPT-2
\end{itemize}
  on SST-2 with ablations over window size, RFF feature count, and gate bias.
\begin{itemize}
\item \textbf{Full reproducibility.} Code, checkpoints, and evaluation scripts released
\end{itemize}
  under {{LICENSE}}.

\subsection{1.5 Paper Organization}

Section 2 reviews related work. Section 3 provides background on local and 
linear attention. Section 4 presents NHA. Section 5 describes integration 
with four architectures. Section 6 details the experimental setup. Section 
7 reports results. Section 8 analyzes gate behavior. Section 9 discusses 
practical guidance. Section 10 lists limitations. Section 11 concludes.



\subsection{Notation}
